\chapter{Methodology}
\lhead{CHAPTER 3. METHODOLOGY}
\label{chap:methodology}

\section{Overview of the Proposed System}
The proposed system combines a Spark--YARN inspired cluster simulator with a distributional deep reinforcement learning scheduler. The pipeline begins with workload ingestion, then constructs the scheduler state, performs policy-based action selection, executes allocations, collects rewards, and updates policy parameters in iterative training episodes \cite{sutton2018rl,bellemare2017distributional}. Conceptually, this is similar to training an operations manager in a virtual factory before deploying them to a real plant: safe experimentation happens first, and decisions improve over repeated feedback.

The complete architecture is referenced in Figure~\ref{fig:system_architecture}. It includes workload generation, cluster state observation, C51 policy network, extension modules (reward optimization, pricing, batching, transfer), and metric logging.

\begin{figure}[H]
\centering
\fbox{\parbox[c][5cm][c]{0.92\textwidth}{\centering Placeholder for system block diagram}}
\caption{System architecture of the proposed distributional deep RL-based Spark job scheduler.}
\label{fig:system_architecture}
\end{figure}

\section{System Components}
\subsection{Apache Spark + YARN Simulation Environment}
The simulation environment models cluster behavior using VM capacities, job resource demands, waiting queues, and execution transitions. Spark jobs are represented as resource-demand vectors (CPU, memory, executor demand, duration), while YARN-like allocation logic enforces capacity constraints per VM \cite{vavilapalli2013apache}. At each decision step, the environment checks whether a selected VM can host a selected job; if yes, allocation proceeds, otherwise the agent receives a penalty for invalid or inefficient action.

This simulation-first approach improves reproducibility and allows systematic stress testing across workload patterns without requiring a costly live cloud deployment.

\subsection{Distributional RL Agent (C51)}
The C51 agent observes state $s_t$, selects action $a_t$, receives reward $R_t$, and transitions to next state $s_{t+1}$. The state space includes queue length, VM availability, and current cost indicators. The action space is a mapping decision: assign job $j_i$ to VM $v_k$, or defer if placement is poor.

Let the categorical support be:
\[
\mathcal{Z} = \{z_1, z_2, \ldots, z_N\}
\]
where $N$ is the number of atoms. The network predicts a probability mass function over this support for each action. The distributional Bellman target is projected onto fixed support using operator $\Pi$.

The C51 optimization objective uses cross-entropy loss:
\[
\mathcal{L}(\theta) = -\sum_{i=1}^{N}\left(\Pi\mathcal{T}Z_{\theta^-}(s_t,a_t)\right)_i \log\left(Z_{\theta}(s_t,a_t)\right)_i
\]
where $\theta$ are online network parameters and $\theta^-$ are target-network parameters \cite{bellemare2017distributional}. In plain English, the model learns to make its predicted reward-distribution match a more accurate target-distribution computed from observed outcomes.

\subsection{Extension 1: BSS Reward Optimization}
The baseline reward in many schedulers tends to over-prioritize either time or cost. To reduce this imbalance, the reward is redesigned as:
\[
R_t = -\alpha\,\text{makespan}_t - \beta\,\text{waste}_t + \gamma\,\text{throughput}_t
\]
where makespan$_t$ is cumulative completion delay at step $t$, waste$_t$ measures under-utilized allocated resources, and throughput$_t$ captures useful completed work.

Each term has intuitive meaning: the first discourages slow schedules, the second discourages resource waste, and the third rewards productive completion. BSS (Bayesian/Black-box Simplex Search style tuning) is then used to find weight combinations that improve long-horizon return under a fixed training budget.

\subsection{Extension 2: Stochastic VM Pricing}
Traditional deterministic pricing assumes fixed VM costs and can mislead scheduling policy in real clouds. We model VM price for machine $v_k$ at time $t$ as a random variable:
\[
P(v_k,t) \sim \mathcal{N}(\mu_k,\sigma_k^2)
\]
or other bounded stochastic process depending on configuration.

The scheduler incorporates expected and observed price fluctuations into state features and reward penalties. This enables cost-aware action choices, such as delaying non-urgent jobs during temporary price spikes when quality-of-service constraints permit.

\subsection{Extension 3: Job Batching}
Instead of scheduling one job at a time, related jobs are grouped into batches:
\[
B = \{b_1,b_2,\ldots,b_q\}
\]
where each batch may group jobs by similar CPU-memory requirements, deadline class, or arrival locality. Batching reduces repetitive per-job decision overhead and improves throughput by improving queue handling efficiency.

In simple terms, this is like handling customer requests in grouped service windows instead of restarting a full decision process for each individual request.

\subsection{Extension 4: Transfer Learning}
Transfer learning pre-trains an agent on source workload $W_1$ and adapts to target workload $W_2$. Lower network layers (feature extractors) are frozen initially to preserve learned representations, while higher decision layers are fine-tuned.

Training-time comparison typically evaluates:
\begin{itemize}
    \item Scratch training on $W_2$
    \item Transfer-frozen (frozen lower layers)
    \item Transfer-finetune (gradually unfreezing layers)
\end{itemize}

This approach reduces retraining effort when workloads evolve, which is common in production systems.

\section{Algorithmic Flow}
\subsection{Algorithm 1: Distributional RL Scheduling Loop}
The end-to-end scheduling loop is described below:
\begin{enumerate}
    \item Initialize environment, replay buffer, C51 online and target networks.
    \item Reset environment and observe initial state $s_0$.
    \item Select action $a_t$ using policy (e.g., $\epsilon$-greedy over distributional value estimates).
    \item Execute action: assign $j_i$ to $v_k$ (or defer), then observe reward $R_t$ and next state $s_{t+1}$.
    \item Store transition $(s_t,a_t,R_t,s_{t+1})$ in replay memory.
    \item Sample mini-batch, compute projected target distribution, and update online network by minimizing C51 loss.
    \item Periodically synchronize target network parameters.
    \item Repeat Steps 3--7 until episode termination.
    \item Log metrics: makespan, utilization, cost, throughput, and convergence indicators.
\end{enumerate}

\subsection{Algorithm 2: Job Batching Pre-Processing}
Batch creation prior to scheduling is performed as follows:
\begin{enumerate}
    \item Collect pending jobs from the waiting queue.
    \item Extract feature vectors (CPU, memory, duration, deadline class).
    \item Compute similarity scores between jobs using normalized resource features.
    \item Group jobs into batches of size $B$ by nearest similarity and arrival constraints.
    \item Order batches by priority policy (deadline-first or arrival-first variants).
    \item Pass batch summary and candidate jobs to scheduling agent as augmented state.
    \item Update batches dynamically as new jobs arrive.
\end{enumerate}

The architecture in Figure~\ref{fig:system_architecture} and these two algorithms jointly define the full methodology used in this project.
